\documentclass[10pt,a4paper]{article}
\usepackage[utf8]{inputenc}
\usepackage[a4paper,left=3cm,right=2cm]{geometry}
\usepackage{amsmath}
\usepackage[thmmarks, amsmath, thref]{ntheorem}
\usepackage{amsfonts}
\usepackage{amssymb}
\usepackage{fancyhdr}
\pagestyle{fancy}
\usepackage{anysize}
\usepackage{graphicx}
\usepackage{float}
\usepackage{hyperref}
\usepackage{multirow}
\usepackage{enumitem}
\newtheorem{theorem}{Theorem}
\theoremsymbol{\ensuremath{\blacksquare}}
\newtheorem*{solution}{Solución:}
\usepackage{listings}
\usepackage{xcolor}

% Margenes y formalidades

\textwidth 6.25in % Margen Derecho y ancho del texto
\textheight 8.5in
\oddsidemargin 0in % Margen Izquierdo
\headheight 0in % Largo del texto
\leftmargin 1in
\parindent 0em % Salto entre la indentación (Sangría)https://www.overleaf.com/4519127669sdpptkhvsbfn
\topmargin -0.5in % Margen Superior
\parskip 0.5ex % Salto entre parrafos
\baselineskip 1pt

%\lhead{\includegraphics[scale=0.2]{logo.png}} % texto izquierda de la cabecera
%\chead{}                                      % texto centro de la cabecera
%\rhead{\includegraphics[scale=0.2]{FIG/logo-mba-2016.png}} % número de página a la derecha

\renewcommand{\headrulewidth}{0.4pt} % grosor de la línea de la cabecera
\renewcommand{\footrulewidth}{0.4pt} % grosor de la línea del pie

\begin{document}



\pagebreak
\vspace*{4mm} 

\begin{center}
\begin{huge}
\textbf{\\Ejercicios Econometría 2}
\end{Large}\\
%{\large 2023-2}\\
\end{center}

\begin{enumerate}

\item  Se tiene los siguentes resultados de un modelo de Regresión Lineal, de las notas finales de un curso, dada caracteristicas del alumno. De los 674 datos se tienen los siguientes resultados de los estimadores:

\begin{center}
    \includegraphics[scale=1]{Ejercicios/sem1-2/eje5.jpg}
\end{center}

\begin{enumerate}
\item Con un nivel de significancia del 5\%, defina cuales son las variables que son significativamente distintas de cero.
\\



\item Haga el intervalo de confianza para el estimador de la variable $attend$, al 95\% de confianza. Concluya.\\


\item Haga el intervalo de confianza para el estimador de la variable $termGPA$, al 95\% de confianza. Concluya.\\


\item Con una significancia del 1\%, testee la hipotyesis de que el estimador de la variable $termGPA$ es igual a 4. Concluya.\\

\end{enumerate}

\item En una muestra de 5 individuos, se encuentra la siguiente matriz inversa:
\[
(X'X)^{-1}=
\begin{pmatrix}
2.7 & 0.3 & -0.65 \\
0.3 & \;\;0.325 & \;\;-0.35 \\
-0.65 & \;\;.0.35 & \;\;0.425 \\
\end{pmatrix}
\]
Se sabe que la varianza del error es de $\sigma^2=3.14$ y los valores de los betas son:
\[
\beta=
\begin{pmatrix}
16.8 \\
2.2 \\
1.4 \\
\end{pmatrix}
\]

El investigador quiere testear la hipotesis de que $2\beta_1=3\beta_2$. ¿Existe evidencia estadística para decir que $2\beta_1=3\beta_2$?



\end{enumerate}
\end{enumerate}




\section{Formulario}


\begin{itemize}

\item  Estadistico t
$$t_{est}=\frac{\hat{\beta_j}-a_j}{s.e.(\hat\beta_j)}$$

    \end{itemize}



\end{document}